\corrigesChapitre{1}{Suites et modèles discrets}
\begin{corrige}{1}{Reconnaître la règle}
\begin{questions}
\item $u_0=5,u_1=8,u_2=11,u_3=14$. La différence vaut $3$.
\item La raison vaut $1/2$ et $v_4=16(1/2)^4=1$.
\item Une suite géométrique de raison $1{,}04$, si le taux est supposé constant.
\end{questions}
\end{corrige}
\begin{corrige}{2}{Indices et sommes}
\begin{questions}
\item Il y a $13$ termes, car les deux bornes sont incluses.
\item $(1-2^5)/(1-2)=31$.
\item Le dernier terme est $u_9=21$ ; la somme vaut $10(3+21)/2=120$.
\end{questions}
\end{corrige}
\begin{corrige}{3}{Limites de référence}
\begin{questions}
\item La première limite vaut $0$ car $0<0{,}6<1$ ; la seconde vaut $+\infty$.
\item $7$ ; puis $(5+2/n)/(2+3/n)\to5/2$.
\item La suite $2-1/(n+1)$ est toujours inférieure à $2$ et tend vers $2$.
\end{questions}
\end{corrige}
\begin{corrige}{4}{Un équilibre}
\begin{questions}
\item $u_1=5$ puis $u_2=5{,}5$.
\item $\ell=0{,}5\ell+3$ donne $\ell=6$.
\item Avec $v_n=u_n-6$, on obtient $v_{n+1}=0{,}5v_n$, $v_0=-2$ ; donc la formule annoncée.
\end{questions}
\end{corrige}
\begin{corrige}{5}{Épargne périodique}
\begin{questions}
\item $u_0=1\,000$, $u_{n+1}=1{,}02u_n+200$.
\item L’équilibre est $-10\,000$ ; $u_n+10\,000$ est géométrique de raison $1{,}02$.
\item $u_2=1\,444{,}4$ euros. La limite est $+\infty$ dans ce modèle à taux et versements constants.
\end{questions}
\end{corrige}
\begin{corrige}{6}{Amortissement}
\begin{questions}
\item $d_0=5\,000$, $d_{n+1}=1{,}05d_n-600$.
\item L’équilibre vaut $12\,000$ ; $d_n=12\,000-7\,000(1{,}05)^n$.
\item Une dette négative ne représente pas la situation annoncée. À la dernière échéance on ajuste le remboursement au montant effectivement dû.
\end{questions}
\end{corrige}
\begin{corrige}{7}{Une somme infinie}
\begin{questions}
\item $S_n=9(1-(1/3)^{n+1})$.
\item $9$, puisque $(1/3)^{n+1}\to0$.
\item $9-S_n=3^{1-n}$. Pour $n=5$ on a $1/81>0{,}01$ ; pour $n=6$, $1/243<0{,}01$. Le premier rang est $6$.
\end{questions}
\end{corrige}
\begin{corrige}{8}{Gendarmes}
\begin{questions}
\item Les deux bornes tendent vers $3$ : le théorème des gendarmes donne $u_n\to3$.
\item $|u_n-3|\leq2/n$, donc $n\geq200$ suffit.
\item Non. Par exemple $u_n=3+(-1)^n/n$ satisfait ces bornes sans être monotone.
\end{questions}
\end{corrige}
\begin{corrige}{9}{Construire et conjecturer}
\begin{questions}
\item $1{,}414$, $1{,}848$, $1{,}961$.
\item Depuis $u_n$ sur l’axe horizontal, monter à la courbe puis aller horizontalement à la diagonale pour reporter la valeur suivante.
\item $\ell=\sqrt{2+\ell}$ implique $\ell^2-\ell-2=0$, donc $\ell=2$ puisque $\ell\geq0$. Cela identifie une limite éventuelle, sans démontrer son existence.
\end{questions}
\end{corrige}
\begin{corrige}{10}{Changer le pas de temps}
\begin{questions}
\item $q^2=1{,}21$ donne $q=1{,}1$.
\item $P_n=100(1{,}1)^n$.
\item Le taux de $21\,\%$ sur deux ans ne se partage pas par addition ; $1{,}105^2=1{,}221025\ne1{,}21$.
\end{questions}
\end{corrige}
\begin{corrige}{11}{Modèle logistique discret}
\begin{questions}
\item $u_1=28$ et $u_2=38{,}08$.
\item $0{,}5\ell(1-\ell/100)=0$ donne $0$ et $100$.
\item Le facteur vaut $1{,}5-0{,}005u_n$ : il diminue lorsque la population augmente. Près de $100$, l’accroissement devient faible ; le modèle géométrique aurait un facteur constant.
\end{questions}
\end{corrige}
\begin{corrige}{12}{Une preuve plutôt qu’un graphique}
\begin{questions}
\item $(0{,}8)^n>0$ donne la borne. $u_{n+1}-u_n=10(0{,}8)^n>0$ donne la croissance.
\item La suite tend vers $150$ ; pour tout $s<150$, elle finit par dépasser $s$. La boucle teste chaque rang dans l’ordre.
\item $0{,}8^{10}\simeq0{,}107374>0{,}1$, mais $0{,}8^{11}\simeq0{,}085899<0{,}1$ : le premier rang est $11$.
\end{questions}
\end{corrige}
